\section{Discussion}
\label{sec:discussion}

\paragraph{Limitations} This is a controlled mechanism study. Both task families are simulated: probes return discrete outcomes from a frozen outcome structure, not log lines, so the translation from raw telemetry to a discrete observation, and what an attacker can do to it, are not measured. The second family removes our own choice of causes, but its outcomes are an LLM's annotation rather than observed telemetry, and replacing one indistinguishable rule conditions it on causes the catalogue can tell apart. Both families are small (24 tasks from eight causes, and 12 rules) and nearly deterministic, development and evaluation episodes share a generator, and our intervals generalize at most to new episodes of these tasks, not to new causes. The candidate causes are given and mutually exclusive, \textit{other} is a residual with no development cases, and the environment announces when the incident changes; new causes, concurrent causes, and undetected change are outside the study. In drifting tasks the switch happens after the second probe, so an agent's pace decides whether it meets the drift. Our verifier is strict, which makes completion conservative. The adversarial tests use four injection texts and two forgeries in one environment, and an attacker who forges two independent source groups would defeat the corroboration rule too. The v1.1 LLM results cover the two small models on a different GPU and serving stack; its primary endpoint on Qwen2.5-72B has not yet run. Public data could not provide a confirmation on real telemetry (\S\ref{sec:realdata}). Finally, we made errors of our own, from a misreported interval and lost journals to reading ExCyTIn test questions during a feasibility check; Appendix~\ref{app:errata} documents each of them, and none affects a confirmed primary endpoint.

\paragraph{Security Implications} The results separate two kinds of attacker influence. Text that reaches the model is contained by the finish guard: an instruction in a log field changed no guarded verdict, whereas every model that decided its own verdict followed it, and removing the raw text alone blocked the dismissal only by leaving most cases unresolved. The injection still costs analyst time, because rejected dismissals end in abstentions and escalations. Evidence that reaches the ledger is not contained: the controller trusts its probes, and acting on the most informative observation first makes a single forged source decisive. Three design rules follow. The guard should always be on, because an unguarded small model overrides the ledger (\S\ref{sec:external}, \S\ref{sec:adversarial}). A benign verdict should require corroboration from distinct source groups, since two probes can read one poisoned feed; which sources are independent is a property of the deployment. And a benign verdict should remain a recommendation that an analyst confirms, so that the residual attack is a delayed response rather than an unattended closure.

\paragraph{Where the Model Is Still Needed} With structured observations, the controller accounts for most of the measured gain. The earlier advantage of strong models that decided when to stop came from waiting for a confirming signature, and a confirmation rule written into the controller reproduces it without an LLM (\S\ref{sec:fair}). The incremental value of an LLM over a controller with an explicit confirmation rule therefore remains to be established; our pre-registered test of it on Qwen2.5-72B is pending. The more plausible role of the model lies outside what we measured: producing the discrete observations from raw logs, handling incomplete descriptions, and noticing conflicting evidence. A real probe returns log lines, not the label \textit{high entropy}. The OTRF case (Appendix~\ref{app:case}) shows what that reading involves: recognizing that a task's creator is a person rather than a machine account, or that a logon came over the network.

\paragraph{Deployment} We recommend \hesp{} with EIG/cost selection, the finish guard, the confirmation stop, and source-group corroboration for benign verdicts, with verdicts handed to an analyst together with their cited evidence and journal. The prediction tables are the main deployment cost. Ours were counted from development episodes that ran every probe on cases of known cause; resolved tickets record only the probes that an analyst chose to run, so tables estimated from them would face selection bias and missing counterfactual outcomes, and would need replayed incidents, expert estimates, or an LLM prior corrected by data. The confirmation rule adds one more requirement: at least one probe must single out each cause, and when the tables blur two causes it escalates rather than guessing, while a cause missing from the tables is the error that produces confident wrong verdicts (\S\ref{sec:mismatch}).
